\documentclass[11pt,a4paper]{article}
\usepackage[margin=2.5cm]{geometry}
\usepackage[T1]{fontenc}
\usepackage{amsmath,amssymb}
\usepackage{booktabs}
\usepackage{url}
\usepackage{hyperref}
\usepackage{listings}
\usepackage{xcolor}
\lstset{basicstyle=\ttfamily\small,breaklines=true,frame=single,backgroundcolor=\color{gray!8},
  literate={├}{{\texttt{|}}}1 {└}{{\texttt{`}}}1 {─}{{\texttt{-}}}1 {│}{{\texttt{|}}}1}
\newcommand{\Ahat}{\widehat{A}}
\newcommand{\dest}[2]{d_{\mathrm{est}}(#1,#2)}

\title{Reconstruction of the duplication history tree: work log}
\author{Atheer A. Matroud\\ \small (documented with Claude)}
\date{18 November 2026}

\begin{document}
\maketitle

\section{Session 1: small example (18 November 2026)}

\subsection{Objective}
Given $k$ homologous tandem repeat sequences, written as TR maps $D_1,\ldots,D_k$, the aim is to
reconstruct their duplication history tree: the tree whose leaves are the $D_i$, whose internal
nodes are reconstructed ancestors, and whose edges carry the duplications, deletions and
substitutions that turn each ancestor into its child, so that every motif copy of every
descendant can be followed back to a motif copy of the root.

\subsection{Decisions taken in this session}
\begin{enumerate}
\item The tree across the homologous sequences and the history of the motif copies inside each
  sequence are reconstructed together, as one object called the duplication history tree.
\item The method extends the \texttt{tr\_aligner} package (asymmetric alignment of Matroud,
  Tuffley and Hendy 2021, and the ancestor reconstruction of the MPE manuscript).
\item The work starts with a small example that can be traced by hand. Once it is validated, a
  simulation follows in which homologous sequences are derived from a five level tree, each
  ancestor being an internal node, with several tree shapes.
\item The small example is a balanced tree with four descendants.
\item Two rules for building the tree are compared: joining the closest pair, and neighbour
  joining.
\item The reconstruction stops at the closest common ancestor of all the sequences. The root is
  not reduced to a single motif copy.
\item Real data, after the simulations: the taro NTR and MSY1 from the Hallast assemblies.
\end{enumerate}

\subsection{The example}
Six variants of length 6 are used (those of the thesis, Chapter~1.3):
\[
a=\texttt{CATGGT},\; b=\texttt{TATGGA},\; c=\texttt{CACGCT},\; d=\texttt{TATGGT},\;
e=\texttt{CAAGGT},\; f=\texttt{CACGGT},
\]
with $H(a,d)=1$ and $H(c,f)=1$. The true history has one event on each edge:
\begin{lstlisting}
R = a b c
  X = a b b c          R2   duplication of b
    D1 = a b b b c     R2   duplication of b
    D2 = a b b f       S4   c -> f
  Y = a b c c          R3   duplication of c
    D3 = a b c c c     R3   duplication of c
    D4 = d b c c       S1   a -> d
\end{lstlisting}
Only $D_1,\ldots,D_4$ are given to the method. Under the edit-weight scoring $\sigma_w$ a
$1$--duplication costs $10$ and a substitution between variants at Hamming distance 1 costs $50$.

\subsection{Method}
\begin{description}
\item[Step 1, distances.] For two TR maps $X$ and $Y$ the ancestor $\Ahat=\Ahat(X,Y)$ is
  reconstructed by the procedure of the MPE manuscript, and
  $d(X,Y)=\dest{\Ahat}{X}+\dest{\Ahat}{Y}$ under the anchored edit-weight scoring.
\item[Step 2, tree.] \emph{Closest pair}: the pair with the smallest $d$ is joined and replaced
  by its reconstructed ancestor, the distances from that ancestor are recalculated, and the step
  is repeated until one sequence remains. \emph{Neighbour joining}: the topology is built once
  from the distances between the descendants, the tree is rooted on the last edge joined, and
  the ancestors are then reconstructed from the leaves upwards.
\item[Step 3, events.] Each ancestor is aligned against each of its children with the anchored
  asymmetric alignment under $\sigma_w$, and the events on the edge are read from the alignment.
\item[Step 4, motif copies.] Each motif copy of a child is assigned the copy of its parent
  against which it was aligned. Composing the edges gives, for every copy of every descendant,
  the copy of the root from which it descends.
\end{description}

\subsection{Results}
\paragraph{Distances.} Table~\ref{tab:dist} gives the six distances and the pair ancestors. Each
distance was checked by hand, and each equals the weight of the true path between the two
descendants in the true tree (for example $D_1$ to $D_3$ is four $1$--duplications, $40$;
$D_1$ to $D_2$ is one duplication and one substitution, $60$).

\begin{table}[h]
\centering
\caption{Pairwise distances through the reconstructed ancestor, anchored edit-weight scoring.}
\label{tab:dist}
\begin{tabular}{llrrr}
\hline\hline
Pair & $\Ahat$ & $\dest{\Ahat}{X}$ & $\dest{\Ahat}{Y}$ & $d(X,Y)$ \\ \hline
$D_1,D_2$ & $a\,b\,b\,f$ & 60 & 0 & 60 \\
$D_1,D_3$ & $a\,b\,c$ & 20 & 20 & 40 \\
$D_1,D_4$ & $d\,b\,c$ & 70 & 10 & 80 \\
$D_2,D_3$ & $a\,b\,c$ & 60 & 20 & 80 \\
$D_2,D_4$ & $d\,b\,c$ & 110 & 10 & 120 \\
$D_3,D_4$ & $a\,b\,c\,c$ & 10 & 50 & 60 \\ \hline\hline
\end{tabular}
\end{table}

\paragraph{Closest pair.} The first join is $(D_1,D_3)$ at distance $40$, which is not a pair of
siblings. The resulting tree is $(D_4,(D_2,(D_1,D_3)))$ with root $d\,b\,c$: wrong topology and
wrong root.
\begin{lstlisting}
A3 = d b c
|- D4 = d b c c            [R3; d_est = 10]
`- A2 = a b f              [S1 S3; d_est = 100]
   |- D2 = a b b f         [R2; d_est = 10]
   `- A1 = a b c           [S3; d_est = 50]
      |- D1 = a b b b c    [R2 R2; d_est = 20]
      `- D3 = a b c c c    [R3 R3; d_est = 20]
\end{lstlisting}

\paragraph{Neighbour joining.} The first join is $(D_1,D_2)$ (the criterion $Q$ is $-320$ for
$(D_1,D_2)$ and $(D_3,D_4)$ and $-280$ for the other four pairs). The tree is
$((D_1,D_2),(D_3,D_4))$, the true topology, and the root is $a\,b\,c$, the true root.
\begin{lstlisting}
A3 = a b c
|- A1 = a b b f            [R2 S3; d_est = 60]
|  |- D1 = a b b b c       [R2 S4; d_est = 60]
|  `- D2 = a b b f         [no events; d_est = 0]
`- A2 = a b c c            [R3; d_est = 10]
   |- D3 = a b c c c       [R3; d_est = 10]
   `- D4 = d b c c         [S1; d_est = 50]
\end{lstlisting}
Motif copies traced to the root $a\,b\,c$ (positions 1, 2, 3):
\begin{center}
\begin{tabular}{lll}
\hline\hline
Descendant & TR map & Root copy of each motif copy \\ \hline
$D_1$ & $a\,b\,b\,b\,c$ & 1, 2, 2, 2, 3 \\
$D_2$ & $a\,b\,b\,f$ & 1, 2, 2, 3 \\
$D_3$ & $a\,b\,c\,c\,c$ & 1, 2, 3, 3, 3 \\
$D_4$ & $d\,b\,c\,c$ & 1, 2, 3, 3 \\ \hline\hline
\end{tabular}
\end{center}
All four agree with the true history.

\subsection{Findings}
\begin{enumerate}
\item \textbf{The closest pair rule fails on this example, and the reason is not the distance.}
  The distances are exact. The true tree is not clock-like: an edge carrying a substitution
  weighs $50$ and an edge carrying a duplication weighs $10$, so two descendants in different
  halves of the tree, separated by four duplications ($40$), are closer than two siblings
  separated by a duplication and a substitution ($60$). Joining the closest pair assumes that
  siblings are the closest pair, which does not hold here. Neighbour joining needs only that the
  distances add along the tree, which they do.
\item \textbf{An internal ancestor with two children cannot decide a substitution.} The true
  $X$ is $a\,b\,b\,c$; the reconstructed $A_1$ is $a\,b\,b\,f$. With $D_1$ ending in $c$ and
  $D_2$ in $f$, the two children alone give no reason to prefer one. As a result the
  substitution $c\to f$ is placed on the wrong edges ($A_3\to A_1$ and back again on
  $A_1\to D_1$), and the total weight of the tree is $190$ instead of the true $140$. The other
  half of the tree ($A_2 = a\,b\,c\,c$) and the root are exact. The parent of $A_1$ holds the
  missing information: a vote at $A_1$ among its two children and its parent would return $c$.
\item \textbf{The position of the root under neighbour joining is a convention.} Neighbour
  joining gives an unrooted tree. Rooting on the last edge joined is correct for the balanced
  example, but with three nodes left the criterion $Q$ is tied ($-160$ for all three pairs), and
  for other tree shapes this convention will not, in general, find the root.
\item The tracing of motif copies to the root was correct for all four descendants, even with
  the wrong variant in $A_1$.
\end{enumerate}

\subsection{Open points for the next session}
\begin{enumerate}
\item Whether to add a correction pass in which each internal ancestor is revised by a vote among
  its two children and its parent (Finding 2).
\item How the root is to be placed on the neighbour-joining tree (Finding 3).
\item Whether the closest pair rule is dropped, or kept for comparison in the simulations.
\item The design of the five level simulation: tree shapes, events per edge, ancestor length.
\end{enumerate}

\subsection{Files}
All in the folder \texttt{Tandem Repeats/DHT}:
\begin{itemize}
\item \texttt{dht/tree.py}: distances, the two tree rules, edge events, tracing of motif copies,
  Newick output.
\item \texttt{examples/small\_example.py} and \texttt{examples/small\_example\_output.txt}: the
  example of this session and its full output, including every alignment.
\item \texttt{tests/test\_small\_example.py}: checks of the results above.
\item \texttt{tr\_aligner/}: unchanged copy of the alignment and ancestor reconstruction package.
\end{itemize}
To reproduce: \texttt{python examples/small\_example.py}.

\end{document}
